\PassOptionsToPackage{unicode=true}{hyperref} % options for packages loaded elsewhere
\PassOptionsToPackage{hyphens}{url}
%
\documentclass[10pt,xcolor=table,color={dvipsnames,usenames},ignorenonframetext,usepdftitle=false,french]{beamer}
\setbeamertemplate{caption}[numbered]
\setbeamertemplate{caption label separator}{: }
\setbeamercolor{caption name}{fg=normal text.fg}
\beamertemplatenavigationsymbolsempty
\usepackage{caption}
\captionsetup{skip=0pt,belowskip=0pt}
%\setlength\abovecaptionskip{-15pt}
\usepackage{lmodern}
\usepackage{amssymb,amsmath,mathtools,multirow}
\usepackage{float,hhline}
\usepackage{tikz}
\usepackage{ifxetex,ifluatex}
\usepackage{fixltx2e} % provides \textsubscript
\ifnum 0\ifxetex 1\fi\ifluatex 1\fi=0 % if pdftex
  \usepackage[T1]{fontenc}
  \usepackage[utf8]{inputenc}
  \usepackage{textcomp} % provides euro and other symbols
\else % if luatex or xelatex
  \usepackage{unicode-math}
  \defaultfontfeatures{Ligatures=TeX,Scale=MatchLowercase}
\fi
\usetheme[coding=utf8,language=french,
,titlepagelogo=img/logobeamer
]{TorinoTh}
% use upquote if available, for straight quotes in verbatim environments
\IfFileExists{upquote.sty}{\usepackage{upquote}}{}
% use microtype if available
\IfFileExists{microtype.sty}{%
\usepackage[]{microtype}
\UseMicrotypeSet[protrusion]{basicmath} % disable protrusion for tt fonts
}{}
\IfFileExists{parskip.sty}{%
\usepackage{parskip}
}{% else
\setlength{\parindent}{0pt}
\setlength{\parskip}{6pt plus 2pt minus 1pt}
}
\usepackage{hyperref}
\hypersetup{
            pdfauthor={Alain QUARTIER-LA-TENTE},
            pdfborder={0 0 0},
            breaklinks=true}
\urlstyle{same}  % don't use monospace font for urls
\newif\ifbibliography
\usepackage{longtable,booktabs}
\usepackage{caption}
% These lines are needed to make table captions work with longtable:
\makeatletter
\def\fnum@table{\tablename~\thetable}
\makeatother
\usepackage{graphicx,grffile}
\makeatletter

\newsavebox\pandoc@box


\newcommand*\pandocbounded[1]{% scales image to fit in text height/width


  \sbox\pandoc@box{#1}%


  \Gscale@div\@tempa{\textheight}{\dimexpr\ht\pandoc@box+\dp\pandoc@box\relax}%


  \Gscale@div\@tempb{\linewidth}{\wd\pandoc@box}%


  \ifdim\@tempb\p@<\@tempa\p@\let\@tempa\@tempb\fi% select the smaller of both


  \ifdim\@tempa\p@<\p@\scalebox{\@tempa}{\usebox\pandoc@box}%


  \else\usebox{\pandoc@box}%


  \fi%


}
% Prevent slide breaks in the middle of a paragraph:
\widowpenalties 1 10000
\raggedbottom
\AtBeginPart{
  \let\insertpartnumber\relax
  \let\partname\relax
  \frame{\partpage}
}
% \AtBeginSection{
%   \ifbibliography
%   \else
%     \begin{frame}{Sommaire}
%     \tableofcontents[currentsection, hideothersubsections]
%     \end{frame}
%   \fi
% }
\setlength{\emergencystretch}{3em}  % prevent overfull lines
\providecommand{\tightlist}{%
  %\setlength{\itemsep}{0pt}
  \setlength{\parskip}{0pt}
  }
\setcounter{secnumdepth}{0}

% set default figure placement to htbp
\makeatletter
\def\fps@figure{htbp}
\makeatother

\usepackage{dsfont}
\usepackage{stmaryrd}
\usepackage[normalem]{ulem}
\usepackage{fontawesome5}
\usepackage{tikz,pgfplots}
\pgfplotsset{compat=1.17}
\pgfplotsset{samples=100}
\usepackage{animate}
 \usepackage{booktabs}

\usepackage{colortbl}

\DeclareMathOperator{\argmin}{argmin}
\DeclareMathOperator{\med}{med}


% blocks
\usepackage{environ}
\usepackage[tikz]{bclogo}

\tikzstyle{titlestyle} =[draw=black!80,fill=black!20, text=black,
 right=10pt, rounded corners]
\mdfdefinestyle{symmaryboxstyle}{
	linecolor=black!80, backgroundcolor = black!5,
	skipabove=\baselineskip, innertopmargin=\baselineskip,
	innerbottommargin=\baselineskip,
	userdefinedwidth=\textwidth,
	middlelinewidth=1.2pt, roundcorner=5pt,
	skipabove={\dimexpr0.5\baselineskip+\topskip\relax},
	frametitleaboveskip=\dimexpr-\ht\strutbox\relax,
	innerlinewidth=0pt,
}
\NewEnviron{summary}{%
\begin{mdframed}[style=symmaryboxstyle]
\vspace{-0.5em}
\BODY
\end{mdframed}
}
\makeatletter
% Open `\noalign` and check for overlay specification:
\def\rowcolor{\noalign{\ifnum0=`}\fi\bmr@rowcolor}
\newcommand<>{\bmr@rowcolor}{%
    \alt#1%
        {\global\let\CT@do@color\CT@@do@color\@ifnextchar[\CT@rowa\CT@rowb}% Rest of original `\rowcolor`
        {\ifnum0=`{\fi}\@gooble@rowcolor}% End `\noalign` and gobble all arguments of `\rowcolor`.
}
% Gobble all normal arguments of `\rowcolor`:
\newcommand{\@gooble@rowcolor}[2][]{\@gooble@rowcolor@}
\newcommand{\@gooble@rowcolor@}[1][]{\@gooble@rowcolor@@}
\newcommand{\@gooble@rowcolor@@}[1][]{\ignorespaces}

\newcommand{\rowc}[1]{\only<#1>{\\\rowcolor{processblue!40}}}
%\newcommand{\rowc}[1]{{\rowcolor<#1>{processblue!30}}
\newcommand{\cellc}[1]{\only<#1>{\cellcolor{processblue!40}}}
\newcommand{\supsp}[1]{\visible<#1>{\\}}
\usepackage{booktabs}
\usepackage{longtable}
\usepackage{array}
\usepackage{multirow}
\usepackage{wrapfig}
\usepackage{float}
\usepackage{colortbl}
\usepackage{pdflscape}
\usepackage{tabu}
\usepackage{threeparttable}
\usepackage{threeparttablex}
\usepackage[normalem]{ulem}
\usepackage{makecell}
\usepackage{xcolor}
\makeatletter
\@ifpackageloaded{caption}{}{\usepackage{caption}}
\AtBeginDocument{%
\ifdefined\contentsname
  \renewcommand*\contentsname{Table of contents}
\else
  \newcommand\contentsname{Table of contents}
\fi
\ifdefined\listfigurename
  \renewcommand*\listfigurename{List of Figures}
\else
  \newcommand\listfigurename{List of Figures}
\fi
\ifdefined\listtablename
  \renewcommand*\listtablename{List of Tables}
\else
  \newcommand\listtablename{List of Tables}
\fi
\ifdefined\figurename
  \renewcommand*\figurename{Figure}
\else
  \newcommand\figurename{Figure}
\fi
\ifdefined\tablename
  \renewcommand*\tablename{Table}
\else
  \newcommand\tablename{Table}
\fi
}
\@ifpackageloaded{float}{}{\usepackage{float}}
\floatstyle{ruled}
\@ifundefined{c@chapter}{\newfloat{codelisting}{h}{lop}}{\newfloat{codelisting}{h}{lop}[chapter]}
\floatname{codelisting}{Listing}
\newcommand*\listoflistings{\listof{codelisting}{List of Listings}}
\makeatother
\makeatletter
\makeatother
\makeatletter
\@ifpackageloaded{caption}{}{\usepackage{caption}}
\@ifpackageloaded{subcaption}{}{\usepackage{subcaption}}
\makeatother

\title{Estimation en temps réel de la tendance-cycle à l'aide de
moyennes mobiles : apports dans l'analyse conjoncturelle}
\ateneo{Soutenance doctorat}
\author{Alain QUARTIER-LA-TENTE}
\date{}


\setrellabel{}

\setcandidatelabel{}

\rel{}
\division{}

\departement{9 décembre 2025}
\makeatletter
\let\@@magyar@captionfix\relax
\makeatother

\DeclareMathOperator{\Cov}{Cov}
\newcommand{\E}[1]{\mathbb{E}\left[ #1 \right]}
\newcommand{\V}[1]{\mathbb{V}\left[ #1 \right]}
\newcommand{\cov}[2]{\Cov\left( #1\,,\,#2 \right)}

\begin{document}
\frame[plain,noframenumbering]{\titlepage}

\section{Introduction}\label{introduction}

\begin{frame}{Introduction}
\phantomsection\label{introduction-1}
Construire des séries CVS-CJO est une pratique courante :

\begin{itemize}
\item
  Permet une analyse en glissement mensuel/trimestriel
\item
  Améliore le diagnostic conjoncturel
\end{itemize}

\begin{center}
\pandocbounded{\includegraphics[keepaspectratio]{img/slides-ffig-ipi-france-1.pdf}}
\end{center}
\end{frame}

\begin{frame}{Hypothèse de travail de la désaisonnalisation}
\phantomsection\label{hypothuxe8se-de-travail-de-la-duxe9saisonnalisation}
Une série brute se décompose en 3 composantes inobservables :

\begin{itemize}
\tightlist
\item
  \textbf{effets saisonniers} (y.c. calendrier) : variations régulières
  le même mois chaque année
\end{itemize}

\pause

\begin{itemize}
\tightlist
\item
  \textbf{tendance-cycle} ou « tendance de court terme » : combinaison
  de la tendance (évolutions de long terme) et du cycle (évolutions de
  court terme autour de la tendance)
\end{itemize}

\pause

\begin{itemize}
\tightlist
\item
  \textbf{irrégulier} : fluctuations inattendues (chocs économiques non
  prévus, bruit\ldots)
\end{itemize}

\pause

\faArrowCircleRight{} La série CVS-CJO contient donc l'effet de la
tendance-cycle \textbf{ET} de l'irrégulier

\pause

Irrégulier pertinent pour l'analyse conjoncturel\ldots{} Mais peut
perturber l'analyse de la tendance sous-jacente et les points de
retournement
\end{frame}

\begin{frame}{Et si l'irrégulier est important ?}
\phantomsection\label{et-si-lirruxe9gulier-est-important}
\begin{center}
\pandocbounded{\includegraphics[keepaspectratio]{img/slides-ffig-ipi-france-sa-1.pdf}}
\end{center}

\faArrowCircleRight{} Dans certains cas, pertinent d'étudier
\textbf{directement} la tendance-cycle
\end{frame}

\begin{frame}{Comment estimer la tendance-cycle ?}
\phantomsection\label{comment-estimer-la-tendance-cycle}
Les deux principaux instituts publiant régulièrement la tendance-cycle
(Statistique Canada et Australian Bureau of Statistics) appliquent une
moyenne mobile (MM) symétrique de \textbf{13 termes} sur la série
désaisonnalisée \textbf{mensuelle} \[
TC_t= \sum_{i=-6}^{+6}\theta_i SA_{t+i}
\] Cohérent avec X-13-ARIMA-SEATS utilisée pour la CVS-CJO

\pause

Les estimations en \textbf{temps réel} s'appuient nécessairement sur des
MM asymétriques qui vont créer des révisions et un déphasage dans la
détection des points de retournement

\pause

Objectifs de la thèse :

\begin{enumerate}
\item
  Étudier des méthodes d'estimation en temps réel de la tendance-cycle
  par des MM \textbf{asymétriques} : revue de l'existant (art. 1) et
  propositions d'extensions (art. 2 et 3)
\item
  Décrire pourquoi publier la tendance-cycle et comment mettre en place
  une production (art. 4)
\end{enumerate}
\end{frame}

\section{Estimation de la tendance-cycle et formule générale de
construction des
MM}\label{estimation-de-la-tendance-cycle-et-formule-guxe9nuxe9rale-de-construction-des-mm}

\begin{frame}[noframenumbering]{Sommaire}
\phantomsection\label{sommaire}
\tableofcontents[currentsection, hideothersubsections]
\end{frame}

\subsection{Introduction}\label{introduction-2}

\begin{frame}{Introduction}
\phantomsection\label{introduction-3}
Récemment, différentes approches générales de construction de MM ont
émergé :

\begin{itemize}
\item
  Filtres polynomiaux locaux
  (\href{https://doi.org/10.1214/08-AOAS195}{Proietti et Luati 2008})
\item
  Filtres fondés sur les espaces de Hilbert à noyau reproduisant (RKHS,
  \href{https://ideas.repec.org/a/bes/jnlbes/v26y2008p536-545.html}{Dagum
  et Bianconcini 2008})
\item
  Approche Fidelity-Smoothness-Timeliness (FST,
  \href{ec.europa.eu/eurostat/web/products-manuals-and-guidelines/-/KS-GQ-18-001}{Grun-Rehomme,
  Guggemos, et Ladiray 2018})
\end{itemize}

\pause

\begin{itemize}
\tightlist
\item
  Approche Accuracy-Timeliness-Smoothness (ATS,
  \href{https://EconPapers.repec.org/RePEc:eee:intfor:v:35:y:2019:i:3:p:1072-1084}{Wildi
  et McElroy 2019})\\
  \faArrowCircleRight{} approche paramétrique non étudiée empiriquement
  car ne correspond pas à la logique non-paramétrique de
  X-13-ARIMA-SEATS
\end{itemize}

\pause

Objectif de l'article 1 : décrire et comparer \textbf{théoriquement} et
\textbf{empiriquement} ces différentes approches pour l'estimation en
temps réel en utilisant une méthodologie commune
\end{frame}

\subsection{Résultats théoriques et
empiriques}\label{ruxe9sultats-thuxe9oriques-et-empiriques}

\begin{frame}{Comparaison théorique}
\phantomsection\label{comparaison-thuxe9orique}
Développement d'une formulation générale pour la construction des
filtres linéaires permettant :

\begin{itemize}[<+->]
\item
  De mettre en évidence les \textbf{équivalences} entre les méthodes et
  des filtres classiques (ex. Macaulay, Nadaraya-Watson)
\item
  De \textbf{clarifier} les \textbf{différences} entre les méthodes
  utilisant des \textbf{critères} du \textbf{même nom} (ex. smoothness)
\item
  De \textbf{clarifier} les \textbf{différences} entre les méthodes se
  basant sur les \textbf{mêmes fondements théoriques} (ex. : RKHS et
  polynômes locaux)
\end{itemize}
\end{frame}

\begin{frame}{Comparaison empirique}
\phantomsection\label{comparaison-empirique}
\begin{itemize}
\tightlist
\item
  Comparaison sur des séries \textbf{mensuelles} réelles et simulées
  avec une MM symétrique de \textbf{13 termes} sur des critères de
  révision et de détection des points de retournement (déphasage)
\end{itemize}

\pause

\begin{itemize}
\tightlist
\item
  \textbf{Fenêtre fixe} appliquée sur \textbf{plusieurs niveaux de
  variabilité} de l'irrégulier (I/C ratio)
\end{itemize}

\pause

Comparaison empirique montre l'importance de bien définir les critères
de comparaison :

\begin{itemize}
\item
  Déphasage : nombre de mois nécessaires pour détecter un point de
  retournement \textbf{sans révision future}\\
  \faArrowCircleRight{} Il peut y avoir des révisions entre les
  estimations intermédiaires successives (ex. RKHS)
\item
  Révisions : par rapport à la dernière estimation mais aussi entre
  \textbf{deux estimations successives} \faArrowCircleRight{} Met en
  évidence des sous-optimalités lorsque l'on se rapproche de
  l'estimation finale (ex. RKHS)
\end{itemize}
\end{frame}

\subsection{Contribution informatique}\label{contribution-informatique}

\begin{frame}[fragile]{Contribution informatique}
\phantomsection\label{contribution-informatique-1}
Développement du package R \texttt{rjd3filters} permettant :

\begin{itemize}
\item
  De facilement créer et combiner différentes moyennes mobiles (et
  notamment reproduire la méthode de décomposition X-11)
\item
  D'utiliser les 4 grandes approches mentionnées tout en étudiant des
  extensions possibles (ex. articles 2 et 3)
\end{itemize}
\end{frame}

\section{Extensions des MM de Henderson et
Musgrave}\label{extensions-des-mm-de-henderson-et-musgrave}

\begin{frame}[noframenumbering]{Sommaire}
\phantomsection\label{sommaire-1}
\tableofcontents[currentsection, hideothersubsections]
\end{frame}

\subsection{Moyennes mobiles de Henderson et de
Musgrave}\label{moyennes-mobiles-de-henderson-et-de-musgrave}

\begin{frame}{Henderson et Musgrave (1/2)}
\phantomsection\label{henderson-et-musgrave-12}
Dans X-13-ARIMA-SEATS, MM de Henderson (symétrique) et Musgrave
(asymétrique) sont utilisées pour estimer la tendance-cycle (comme le
fait ABS)

\pause

MM de Henderson obtenue par analogie avec la régression polynomiale
locale~: \(y_t = TC(t)+\varepsilon_t\) et \(TC(t)=TC_t\) peut être
\textbf{localement} approchée par un polynôme en \(t\) ~: \[
\forall j\in\left\llbracket -h,h\right\rrbracket : y_{t+j}=\underbrace{\sum_{i=0}^{3}\beta_{i,t}j^{i}}_{\simeq TC_{t+j}}+\varepsilon_{t+j}
\]

\pause

Estimation en utilisant les WLS : \(\hat{\beta_t}=(X'KX)^{1}X'Ky_t\) et
\[
\widehat{TC}_{t}=\hat\beta_{0,t}=w'y=\sum_{j=-h}^{h}w_{j}y_{t+j}
\] Avec des poids biens choisis on retrouve la MM \textbf{symétrique} de
Henderson
\end{frame}

\begin{frame}{Henderson et Musgrave (2/2)}
\phantomsection\label{henderson-et-musgrave-22}
MM de Musgrave obtenue par \textbf{compromis biais-variance} \[
\forall j\in\left\llbracket -h,f\right\rrbracket : y_{t+j}=\beta_{0,f}+\delta t+\varepsilon_{t+j}\quad\varepsilon_t\sim \mathcal N(0,\sigma^2).
\] Coefficients obtenus en minimisant l'EQM de révision sous contrainte
de \textbf{préservation sans biais des constantes} (\(\beta_{0,f}\))
mais avec \textbf{biais} dans l'estimation de la \textbf{pente}
(\(\delta\))

\pause

MM dépend du rapport \(|\delta/\sigma|\) défini par l'utilisateur. Si
biais (\(\delta\)) constant dans le temps, tendance est linéaire,
peut-être lié à l'I-C ratio (\(R\)) \[
\left|\frac{\delta}{\sigma}\right|=\frac{2}{R\sqrt{\pi}}.
\]

C'est ce qui fait dans X-13-ARIMA-SEATS en utilisant \(R=3,5\) lorsque
H-13 termes est utilisée (ABS)

\pause

Objectifs articles 2 et 3 : étudier des extensions pour améliorer la
performance en temps réel
\end{frame}

\subsection{Paramétrisation locale des filtres de régression
polynomiale}\label{paramuxe9trisation-locale-des-filtres-de-ruxe9gression-polynomiale}

\begin{frame}{Estimation locale de \(|\delta/\sigma|\)}
\phantomsection\label{estimation-locale-de-deltasigma}
\textbf{Hypothèses} de biais constant et de tendance linéaire souvent
\textbf{irréalistes}, notamment autour des points de retournement !

\pause

Solution : \textbf{estimer localement} \(|\delta/\sigma|\) :

\begin{itemize}
\tightlist
\item
  Pente \(\delta\) estimée localement via une régression polynomiale
  locale (construction de MM pour estimer)\\
\item
  Variance \(\sigma^2\) estimée en utilisant l'ensemble des données
\end{itemize}

\begin{center}
\pandocbounded{\includegraphics[keepaspectratio]{img/slides-ffig-local-icr-1.pdf}}
\end{center}

\pause

Résultats sur des séries mensuelles (simulées et réelles) :
\textbf{réduit} les \textbf{révisions} et le \textbf{délai dans la
détection des points de retournement}

\textbf{Pas d'impact} significatif sur des séries
\textbf{trimestrielles}
\end{frame}

\subsection{Construction de moyennes mobiles robustes aux points
atypiques}\label{construction-de-moyennes-mobiles-robustes-aux-points-atypiques}

\begin{frame}{Impact des points atypiques (1/3)}
\phantomsection\label{impact-des-points-atypiques-13}
\begin{itemize}[<+->]
\item
  Par construction les MM sont \textbf{sensibles aux points atypiques}
\item
  Solution 1 : Correction a priori via un modèle RegARIMA
  (X-13-ARIMA-SEATS et ABS sauf COVID-19) ? Inconvénients : convergence
  peut être longue, estimation dépend de la spécification du modèle et
  de la période d'estimation
\end{itemize}

\begin{center}
\pandocbounded{\includegraphics[keepaspectratio]{img/slides-ffig-ex-simulchoc-1.pdf}}
\end{center}
\end{frame}

\begin{frame}{Impact des points atypiques (2/3)}
\phantomsection\label{impact-des-points-atypiques-23}
\begin{itemize}
\tightlist
\item
  Solution 2 pour les ruptures en niveau (LS), possibilité d'une
  estimation segmentée (Statistique Canada)
\end{itemize}

\begin{center}
\pandocbounded{\includegraphics[keepaspectratio]{img/slides-ffig-ex-simulchoc2-1.pdf}}
\end{center}
\end{frame}

\begin{frame}{Impact des points atypiques (3/3)}
\phantomsection\label{impact-des-points-atypiques-33}
Objectif de l'article : étudier des approches alternatives :

\begin{enumerate}
\tightlist
\item
  Étude de méthodes \textbf{non-linéaires robustes} (médianes mobiles,
  moindres carrés médians\ldots)
\end{enumerate}

\pause

\begin{enumerate}
\setcounter{enumi}{1}
\tightlist
\item
  Exploiter la modélisation polynomiale pour construire des \textbf{MM
  intégrant une information extérieure} (choc ponctuel ou en niveau)
  afin d'améliorer la décomposition tendance-cycle/irrégulier
\end{enumerate}

\begin{center}
\pandocbounded{\includegraphics[keepaspectratio]{img/slides-ffig-ex-simulchoc4-1.pdf}}
\end{center}
\end{frame}

\begin{frame}{Principaux résultats}
\phantomsection\label{principaux-ruxe9sultats}
Méthodes non-linéaires robustes produisent de \textbf{fortes révisions
en temps réel} et ne reproduisent pas les points de retournement
associés aux chocs (ex COVID-19)

\pause

MM «~robustes~» : \textbf{réduisent les révisions} et
\textbf{modélisent} bien les \textbf{points de retournements}

\pause

Problématique de la \textbf{validation} et de l'\textbf{identification}
et des chocs à modéliser :

\begin{itemize}
\item
  Validation par comparaison avec les estimations sans choc et en
  construisant des \textbf{intervalles de confiance}
\item
  \textbf{Pas d'impact} à modéliser un \textbf{choc inexistant}
\item
  Utilisation d'information \emph{ad hoc} ou de modèles externes (ex.
  RegARIMA) pour l'identification des chocs
\item
  Révisions importantes si la nature du choc est mal modélisée
  (association du choc à l'irrégulier ou la tendance-cycle) : en cas
  d'incertitude il vaut mieux ne rien modéliser
\end{itemize}
\end{frame}

\section{Et si l'on publiait la
tendance-cycle~?}\label{et-si-lon-publiait-la-tendance-cycle}

\begin{frame}[noframenumbering]{Sommaire}
\phantomsection\label{sommaire-2}
\tableofcontents[currentsection, hideothersubsections]
\end{frame}

\begin{frame}{Pourquoi et comment publier la TC ? (1/2)}
\phantomsection\label{pourquoi-et-comment-publier-la-tc-12}
\href{https://www.imf.org/external/pubs/ft/qna/pdf/2017/QNAManual2017FRE.pdf}{FMI
(2017)} et
\href{https://ec.europa.eu/eurostat/documents/3859598/19355229/KS-GQ-24-012-EN-N.pdf/be954db6-64f5-c1a2-a0b8-0b4de2a5c707?version=1.0&t=1718182927069}{Eurostat
(2024)} recommandent de publier la TC en plus de la série
CVS-CJO\ldots{} Mais très peu d'instituts publient régulièrement la TC

\pause

Objectifs de l'article :

\begin{enumerate}
\tightlist
\item
  Synthétiser les arguments en faveur de la publication de la TC et les
  recommandations pour le faire
\end{enumerate}

\pause

Trois principales raisons de publier la TC :

\begin{enumerate}
[i.]
\tightlist
\item
  Réduire le risque que les utilisateurs tirent des \textbf{conclusions
  inappropriées} sur la base de mouvements liés à l'irrégulier\\
\item
  Permettre une \textbf{comparaison} appropriée \textbf{dans le temps}
  en réduisant l'impact des évènements ponctuels\\
\item
  Améliorer la \textbf{compréhension} et la détection des points de
  retournement
\end{enumerate}

Recommandations : présenter la TC avec la série CVS-CJO et refléter
l'incertitude des premières estimations
\end{frame}

\begin{frame}[fragile]{Pourquoi et comment publier la TC ? (2/2)}
\phantomsection\label{pourquoi-et-comment-publier-la-tc-22}
\begin{enumerate}
\setcounter{enumi}{1}
\item
  Décrire comment, en pratique, publier la TC grâce au package
  \texttt{publishTC} permettant :

  \begin{itemize}
  \item
    L'estimation simple de la TC avec les méthodes de Statistique
    Canada, ABS et les extensions proposées dans les articles 2 et 3
    (ainsi que combinaison de ces extensions)
  \item
    Différentes visualisations (y.c. intervalles de confiance,
    prévisions implicites\ldots)
  \item
    Recommandations sur le choix de la méthode et critères de
    comparaison
  \item
    Création d'un \emph{template} simplifié de mise en production
    automatique appliqué à
    \href{https://aqlt.github.io/publishTC.wp/}{10 publications de
    l'Insee} (\(\simeq 80\) séries)
  \end{itemize}
\end{enumerate}
\end{frame}

\section{Conclusion}\label{conclusion}

\begin{frame}[fragile]{Conclusion}
\phantomsection\label{conclusion-1}
Apports de la thèse :

\begin{itemize}
\item
  \textbf{Unification des méthodes} de construction des moyennes
  mobiles\\
  \faArrowCircleRight{} développement d'une formulation générale et
  comparaison selon une méthodologie commune
\item
  \textbf{Deux améliorations} des méthodes existantes :

  \begin{itemize}
  \item
    Paramétrisation locale des moyennes mobiles asymétriques
    \faArrowCircleRight{} réduction des révisions et délais dans la
    détection des points de retournements
  \item
    Moyennes mobiles robustes aux chocs \faArrowCircleRight{}
    estimations de la TC plus proches de la réalité économique
  \end{itemize}
\item
  \textbf{Outils open-source} : \texttt{rjd3filters} et
  \texttt{publishTC} \faArrowCircleRight{} facilitent la recherche et la
  mise en production
\end{itemize}

\pause

Exemples de prolongements : travaux sur la sélection de la fenêtre (fixe
ou mobile), intégration des moyennes mobiles «~robustes~» dans
X-11\ldots{}
\end{frame}

\begin{frame}[fragile,noframenumbering]{Merci pour votre attention}
\phantomsection\label{merci-pour-votre-attention}
\footnotesize

\begin{itemize}
\item
  Article 1 : Estimation en temps réel de la tendance-cycle : apport de
  l'utilisation des moyennes mobiles asymétriques, \emph{Document de
  travail Insee}, M2024/01.
  \href{https://inseefrlab.github.io/DT-est-tr-tc/}{inseefrlab.github.io/DT-est-tr-tc}
\item
  Article 2 : Improving Real-Time Trend Estimates Using Local
  Parametrization of Polynomial Regression Filters. \emph{Journal of
  Official Statistics, 40}(4), 685-715.
  \href{https://github.com/AQLT/joffstats2024}{github.com/AQLT/joffstats2024}
\item
  Article 3 : Estimation de la tendance-cycle avec des méthodes robustes
  aux points atypiques.
  \href{https://aqlt.github.io/robustMA/}{aqlt.github.io/robustMA}
\item
  Article 4 : Et si l'on publiait la tendance-cycle~?, \emph{XVèmes
  Journées de Méthodologie Statistique de l'Insee}.
  \href{https://github.com/AQLT/publishTC.wp}{github.com/AQLT/publishTC.wp}
\end{itemize}

Packages \faIcon{r-project} : \texttt{rjd3filters}
(\href{https://github.com/rjdverse/rjd3filters}{github.com/rjdverse/rjd3filters})
et \texttt{publishTC}
(\href{https://github.com/AQLT/publishTC}{github.com/AQLT/publishTC})
\end{frame}

\end{document}
